% 张三慧第四版第二卷第15章《静电场中的电介质》独立解答
% 原教材扫描物理页99--102。全部是未经外部学术复核的候选稿。
\section*{第15章\quad 静电场中的电介质：独立解析}

\paragraph{习题 15.1}
取氯核位置为 \(x=0\)，质子在 \(x=d=0.128\,\mathrm{nm}\)；
氯核电荷 \(+17e\)，氢核 \(+e\)，整分子18个电子
的总负电荷为 \(-18e\)。
若全部负电荷球对称集中于氯核，则总电矩
\[
 \boxed{p_{\rm ion}=ed=(1.60218\times10^{-19})
 (0.128\times10^{-9})
 =2.05\times10^{-29}\,\mathrm{C\,m}}.
\]
设真实负电荷分布重心位于
\(x_{\rm neg}\)，
则沿键轴总电矩
\(p=ed-18ex_{\rm neg}\)（按正向取符号）。
取实测电矩正向与质子位置一致，得到
\[
 \boxed{x_{\rm neg}=
 \frac{d-p_{\rm meas}/e}{18}
 \approx5.93\times10^{-12}\,\mathrm m}
\]
即从氯核朝质子方向约 \(0.00593\,\mathrm{nm}\)。
不要把整个负电荷云的电量误当作只有一个电子；
若只求\emph{有效正负单元电荷}的分离量，
\(p/e\approx0.0212\,\mathrm{nm}\)，
但这不是18电子分布重心距氯核的距离。

\paragraph{习题 15.2}
金属球壳间 \(R_1<r<R_2\) 的球对称自由电荷
只有内壳所带 \(Q\)。
无论介质如何分层，高斯定律
\(\oiint\boldsymbol D\cdot d\boldsymbol S=Q\) 给出
\[
 \boxed{\boldsymbol D(r)=
 \frac{Q}{4\pi r^2}\,\hat{\boldsymbol r}},\qquad
 \boxed{\boldsymbol E(r)=
 \begin{cases}
 \dfrac{Q}{4\pi\varepsilon_0\varepsilon_{r1}r^2}
 \hat{\boldsymbol r},&R_1<r<R,\\[3pt]
 \dfrac{Q}{4\pi\varepsilon_0\varepsilon_{r2}r^2}
 \hat{\boldsymbol r},&R<r<R_2.
 \end{cases}}
\]
在金属内部 \(E=D=0\)；两介质交界处
\(D_r\) 连续，\(E_r\) 随介电常量跳变。
取 \(U_{12}=\varphi(R_1)-\varphi(R_2)\)，
\[
 \boxed{U_{12}=\frac{Q}{4\pi\varepsilon_0}
 \left[\frac1{\varepsilon_{r1}}
 \left(\frac1{R_1}-\frac1R\right)
 +\frac1{\varepsilon_{r2}}
 \left(\frac1R-\frac1{R_2}\right)\right]}
 \approx-8.99\times10^2\,\mathrm V.
\]
贴近内金属球壳的介质\emph{内侧}法线指向
\(-\hat{\boldsymbol r}\)，
故介质束缚面电荷密度
\[
 \boxed{\sigma_b(R_1^+)=
 -\frac{\varepsilon_{r1}-1}{\varepsilon_{r1}}
 \frac{Q}{4\pi R_1^2}
 \approx+9.95\times10^{-6}\,\mathrm{C/m^2}}.
\]
因内壳为负电，其相邻介质表面的束缚电荷为正。

\paragraph{习题 15.3}
设内导体每单位长度带自由电量 \(\lambda\)，
两介质分层交界半径 \(r_0\)，
则在同轴电容器中
\(D_r=\lambda/(2\pi r)\)，
\[
 E_1(r)=\frac{\lambda}{2\pi\varepsilon_0
 \varepsilon_{r1}r},\quad
 E_2(r)=\frac{\lambda}{2\pi\varepsilon_0
 \varepsilon_{r2}r}.
\]
各层最大的场强分别发生在 \(R_1\) 及 \(r_0\)：
\[
 \frac{E_{2,\max}}{E_{1,\max}}
 =\frac{\varepsilon_{r1}R_1}
 {\varepsilon_{r2}r_0}
 =\frac{2R_1}{r_0}>1
\]
因为 \(r_0<R_2<2R_1\)。
所以\textbf{外层第二介质先击穿}。
临界线电荷量
\(\lambda_{\max}=2\pi\varepsilon_0
 \varepsilon_{r2}r_0E_{\max}\)，
\[
 \boxed{U_{\max}=E_{\max}\,r_0
 \left[\frac12\ln\frac{r_0}{R_1}
 +\ln\frac{R_2}{r_0}\right]}.
\]
该式基于题设两介质具有相同击穿场强；
若不同则应比较两者分别的临界电压。

\paragraph{习题 15.4}
无限大板近似下，
金属板自由面电荷密度为
\(\sigma_f=Q/S\)，
介质电位移大小
\(D=|\sigma_f|\)，
介质电场大小
\(E=|\sigma_f|/(\varepsilon_0\varepsilon_r)\)。
极化强度为
\[
 P=\varepsilon_0(\varepsilon_r-1)E
 =\left(1-\frac1{\varepsilon_r}\right)
 |\sigma_f|.
\]
介质相邻金属板面上的束缚电荷
与金属自由电荷异号，
其大小积分为
\[
 \boxed{|Q_b|=PS
 =\left(1-\frac1{\varepsilon_r}\right)|Q|}.
\]
若教材的 \(Q,Q'\) 均表示正的电量大小，
即可写为题给 \(Q'=(1-1/\varepsilon_r)Q\)。

\paragraph{习题 15.5}
无限长细圆柱导线外近似有
\(E(R)=|\lambda|/(2\pi\varepsilon_0R)\)。
设空气击穿场强
\(E_{\max}=3.0\times10^6\,\mathrm{V/m}\)，
因此
\[
 \boxed{|\lambda|_{\max}
 =2\pi\varepsilon_0R E_{\max}}.
\]
代入题给三个半径分别得
\[
 \boxed{\begin{aligned}
 R=1.0\,\mathrm{cm}:&\quad1.67\times10^{-6}\,\mathrm{C/m},\\
 R=1.0\,\mathrm{mm}:&\quad1.67\times10^{-7}\,\mathrm{C/m},\\
 R=0.1\,\mathrm{mm}:&\quad1.67\times10^{-8}\,\mathrm{C/m}.
 \end{aligned}}
\]
真实有限长导线端部曲率更大、表面电场集中，
本结果仅针对忽略端部效应的模型。

\paragraph{习题 15.6}
细胞膜厚
\(d=5.2\times10^{-9}\,\mathrm m\)，
自由表面电荷密度大小
\(\sigma_f=0.52\times10^{-3}\,\mathrm{C/m^2}\)，
相对介电常量 \(\varepsilon_r=6.0\)。
把膜近似为平行板结构，
\[
 \boxed{E=\frac{\sigma_f}
 {\varepsilon_0\varepsilon_r}
 \approx9.79\times10^6\,\mathrm{V/m}},
\quad
 \boxed{\Delta V=Ed\approx0.0509\,\mathrm V}.
\]
外表面为正，膜内电场从外表面指向内表面，
外表面电势高于内表面 \(0.0509\,\mathrm V\)。
忽略膜上离子通道、双电层与非均匀介质响应。

\paragraph{习题 15.7}
板与外加电场 \(E_0\) 的夹角为 \(\theta\)，
故电场垂直于板面的法向分量
\(E_{0,n}=E_0\sin\theta\)。
假设无限大绝缘介质平板，
板界面无自由面电荷，
法向电位移连续
\(\varepsilon_0E_{0,n}
 =\varepsilon_0\varepsilon_r E_{{\rm in},n}\)，
因此
\(E_{{\rm in},n}=E_0\sin\theta/\varepsilon_r\)。
极化法向量
\(P_n=\varepsilon_0(\varepsilon_r-1)E_{{\rm in},n}\)。
两表面束缚面电荷密度分别为
\[
 \boxed{\sigma_b=\pm\varepsilon_0
 \left(1-\frac1{\varepsilon_r}\right)
 E_0\sin\theta}.
\]
正负号由各面的外法线朝向决定；
若把 \(\theta\) 误当成电场与板法线夹角，
会误写成 \(\cos\theta\)。

\paragraph{习题 15.8}
两片空气电容器初距离
\(d_0=0.600\,\mathrm{mm}\)，面积
\(S=50.0\,\mathrm{mm^2}=5.0\times10^{-5}\,\mathrm{m^2}\)。
初电容
\[
 C_0=\frac{\varepsilon_0S}{d_0}
 \approx0.738\,\mathrm{pF}.
\]
触发阈值为增加
\(\Delta C=0.250\,\mathrm{pF}\)。
末距离
\(d_f=\varepsilon_0S/(C_0+\Delta C)\)，
故下压量
\[
 \boxed{\Delta d=d_0-d_f
 =d_0-\frac{\varepsilon_0S}
 {\varepsilon_0S/d_0+\Delta C}
 \approx0.152\,\mathrm{mm}}.
\]
不允许把 \(0.250\,\mathrm{pF}\)
误当作按下后的\emph{总电容}。

\paragraph{习题 15.9}
薄膜厚
\(d=5.0\times10^{-5}\,\mathrm m\)，
\(\varepsilon_r=2.3\)，电容要求
\(C=3.0\times10^{-6}\,\mathrm F\)。
由平行板电容
\(C=\varepsilon_0\varepsilon_rS/d\)，
所需单层有效相对电极面积
\[
 \boxed{S=\frac{Cd}{\varepsilon_0\varepsilon_r}
 \approx7.37\,\mathrm{m^2}}.
\]
这是理想单层薄膜面积；若卷绕成多层电容器，
需要说明叠层及电极连接方式。

\paragraph{习题 15.10}
击穿场强为
\(E_{\rm bd}=3\times10^6\,\mathrm{V/m}\)，
板间电压 \(U=50\,\mathrm{kV}\)。
为不击穿必须有
\(d\ge U/E_{\rm bd}=1/60\,\mathrm m\)。
空气板的面积电容
\(C/S=\varepsilon_0/d\)，
最大值发生在最小允许间距：
\[
 \boxed{\left(\frac CS\right)_{\max}
 =\frac{\varepsilon_0 E_{\rm bd}}{U}
 \approx5.31\times10^{-10}\,\mathrm{F/m^2}}.
\]

\paragraph{习题 15.11}
孤立导体球电容 \(C=4\pi\varepsilon_0R\)。
取 \(R=0.18\,\mathrm m\)，
\[
 \boxed{C\approx2.00\times10^{-11}\,\mathrm F
 =20.0\,\mathrm{pF}},
\quad
 \boxed{Q=CU
 \approx4.00\times10^{-6}\,\mathrm C}.
\]
此结果未计加速器支撑件、相邻电极和环境导体的影响。

\paragraph{习题 15.12}
同轴真空导体构成长度为 \(L\) 的圆柱电容器，
内丝半径
\(a=1.25\times10^{-5}\,\mathrm m\)，
圆筒内半径
\(b=1.25\times10^{-2}\,\mathrm m\)，
长度 \(L=0.100\,\mathrm m\)。
由同轴场
\(E_r=\lambda/(2\pi\varepsilon_0r)\)，
电压为
\(U=\lambda\ln(b/a)/(2\pi\varepsilon_0)\)，
故
\[
 \boxed{C=\frac{2\pi\varepsilon_0 L}{\ln(b/a)}
 =\frac{2\pi\varepsilon_0(0.100)}{\ln1000}
 \approx8.05\times10^{-13}\,\mathrm F
 =0.805\,\mathrm{pF}}.
\]
忽略两端边缘场。

\paragraph{习题 15.13}
当相邻极板交替归属两组电极时，
\(n\) 块板之间共有 \(n-1\) 个相邻间隙。
每个间隙在最大重叠时的电容为
\(C_0=\varepsilon_0S/d\)。
各间隙两端分别连接同一对电极端子，
故电气上\emph{并联}，而不是串联：
\[
 \boxed{C_{\max}=(n-1)C_0
 =\frac{(n-1)\varepsilon_0S}{d}}.
\]
假定所有极板面积与间距相同，并忽略边缘效应。

\paragraph{习题 15.14}
设内部上、下两极板为A、B，
盒体为第三个等势导体C。
除板间直接电容
\(C_{AB}=\varepsilon_0S/d\) 之外，
各板分别和盒体之间有
\(C_{AC}=C_{BC}=\varepsilon_0S/s\)，
其中
\(d=0.50\,\mathrm{mm}\)，
\(s=0.25\,\mathrm{mm}\)。
当盒体\emph{悬浮且净电量为零}，
A—C—B支路相当于两等值电容串联，
与A—B直接电容并联：
\[
 \boxed{C_{\rm floating}
 =\frac{\varepsilon_0S}{d}
 +\frac12\frac{\varepsilon_0S}{s}
 \approx7.08\times10^{-10}\,\mathrm F}.
\]
若A与盒体C短接，
两者相同电势，B对A/C的两个电容并联：
\[
 \boxed{C_{\rm shorted}
 =\varepsilon_0S\left(\frac1d+\frac1s\right)
 \approx1.06\times10^{-9}\,\mathrm F}.
\]
若金属盒接地或有其他外部连接，
有效两端网络的边界条件需重新确定。

\paragraph{习题 15.15}
令板长方向为 \(x\in[0,a]\)，
宽 \(b\)，局部间距
\(h(x)=d+(l/a)x\)；
将板分割为并联的宽度为 \(dx\) 的窄条，
\[
 dC=\frac{\varepsilon_0 b\,dx}{h(x)}.
\]
积分得到
\[
 \boxed{C=\frac{\varepsilon_0ab}{l}
 \ln\left(1+\frac ld\right)}.
\]
在 \(l\ll d\) 时，
\[
 C=\frac{\varepsilon_0ab}{d}
 \left(1-\frac l{2d}+O(l^2/d^2)\right),
\]
恢复平行板近似。
这需要板倾斜角很小且边缘场可忽略。
